
\documentclass{article}
\usepackage{colortbl}
\usepackage{makecell}
\usepackage{multirow}
\usepackage{supertabular}

\begin{document}

\newcounter{utterance}

\centering \large Interaction Transcript for game `cladder', experiment `full\_v1.5\_default', episode 6302 with SLED{-}BSU\_\_Qwen3.5{-}27B{-}sft{-}ep1\_\_prm{-}ep1{-}best\_of\_n.
\vspace{24pt}

{ \footnotesize  \setcounter{utterance}{1}
\setlength{\tabcolsep}{0pt}
\begin{supertabular}{c@{$\;$}|p{.15\linewidth}@{}p{.15\linewidth}p{.15\linewidth}p{.15\linewidth}p{.15\linewidth}p{.15\linewidth}}
    \# & \multicolumn{2}{c}{Player} && \multicolumn{2}{c}{Game Master} \\
    \hline

    \theutterance \stepcounter{utterance}  
    & & & \multicolumn{4}{p{0.6\linewidth}}{
        \cellcolor[rgb]{0.9,0.9,0.9}{
            \makecell[{{p{\linewidth}}}]{
                \texttt{\tiny{[P1$\langle$GM]}}
                \texttt{You are an expert in causal inference. The following question is not a typical commonsense query, but rather a meticulously designed question created by a professor specializing in causal inference, intended to assess the students' mastery of the course content.} \\
\texttt{Imagine a self{-}contained, hypothetical world with only the following conditions, and without any unmentioned factors or causal relationships: Zuph has a direct effect on rixq and xevu. Rixq has a direct effect on xevu. For those who are not zuph and are not rixq, the probability of xevu is 15\%. For those who are not zuph and are rixq, the probability of xevu is 50\%. For those who are zuph and are not rixq, the probability of xevu is 32\%. For those who are zuph and are rixq, the probability of xevu is 85\%. The overall probability of zuph is 84\%. Will rixq increase the chance of xevu?} \\
\texttt{Start your answer with 'Yes' or 'No', followed by additional reasoning or evidence to support your explanation.} \\
            }
        }
    }
     \\ \\

    \theutterance \stepcounter{utterance}  
    & \multicolumn{4}{p{0.6\linewidth}}{
        \cellcolor[rgb]{0.561,0.737,0.561}{
            \makecell[{{p{\linewidth}}}]{
                \texttt{\tiny{[P1$\rangle$GM]}}
                \texttt{No.} \\
\\ 
\texttt{Reasoning:  } \\
\texttt{We can compute the probability of xevu for each group and then compare the overall probability of xevu for those who are rixq versus those who are not rixq, using the given conditional probabilities and the overall prevalence of zuph.} \\
\\ 
\texttt{Let  } \\
\texttt{{-} Z = zuph,  } \\
\texttt{{-} R = rixq,  } \\
\texttt{{-} X = xevu.} \\
\\ 
\texttt{Given:  } \\
\texttt{P(Z) = 0.84, so P(¬Z) = 0.16.} \\
\\ 
\texttt{Conditional probabilities:  } \\
\texttt{{-} P(X | ¬Z, ¬R) = 0.15  } \\
\texttt{{-} P(X | ¬Z, R) = 0.50  } \\
\texttt{{-} P(X | Z, ¬R) = 0.32  } \\
\texttt{{-} P(X | Z, R) = 0.85  } \\
\\ 
\texttt{We do not know P(R) directly, but we can express the overall probability of X for R and ¬R in terms of P(R).} \\
\\ 
\texttt{Let p = P(R). Then:} \\
\\ 
\texttt{P(X | R) = P(X | R, Z)P(Z | R) + P(X | R, ¬Z)P(¬Z | R)  } \\
\texttt{P(X | ¬R) = P(X | ¬R, Z)P(Z | ¬R) + P(X | ¬R, ¬Z)P(¬Z | ¬R)} \\
\\ 
\texttt{But we don’t know P(Z | R) or P(Z | ¬R). However, the question asks whether “rixq will increase the chance of xevu” in this world. In causal terms, we should check whether the *causal effect* of R on X is positive, i.e., whether intervening to set R = 1 (do(R)) increases X compared to do(R = 0), holding Z fixed or averaging over Z.} \\
\\ 
\texttt{Because Z is a cause of both R and X (Zuph has a direct effect on rixq and xevu), R is a *mediator* on the path Z → R → X, but also there is a direct path Z → X. In this simple DAG:} \\
\\ 
\texttt{Z → R  } \\
\texttt{Z → X  } \\
\texttt{R → X} \\
\\ 
\texttt{The causal effect of R on X is the direct effect R → X, which we can identify by comparing X when R = 1 versus R = 0 *within each level of Z* (since Z is a parent of R, conditioning on Z blocks the backdoor path).} \\
\\ 
\texttt{Within Z = 0 (¬Z):  } \\
\texttt{{-} P(X | R=1, Z=0) = 0.50  } \\
\texttt{{-} P(X | R=0, Z=0) = 0.15  } \\
\texttt{So the effect of R on X when Z=0 is 0.50 − 0.15 = +0.35.} \\
\\ 
\texttt{Within Z = 1:  } \\
\texttt{{-} P(X | R=1, Z=1) = 0.85  } \\
\texttt{{-} P(X | R=0, Z=1) = 0.32  } \\
\texttt{So the effect of R on X when Z=1 is 0.85 − 0.32 = +0.53.} \\
\\ 
\texttt{In both strata of Z, R increases X. Therefore, the *causal* effect of rixq on xevu is positive.} \\
\\ 
\texttt{If the question is interpreted as “does being rixq increase the chance of xevu” in the causal sense (i.e., does intervening to make someone rixq increase their probability of xevu), the answer is **Yes**.} \\
\\ 
\texttt{However, the question as phrased—“Will rixq increase the chance of xevu?”—is ambiguous between:} \\
\texttt{{-} Causal effect (do(R)): Yes, as shown above.} \\
\texttt{{-} Observational association P(X | R) vs P(X | ¬R): This depends on the unknown distribution of R within Z, and could be different.} \\
\\ 
\texttt{But in the context of a causal inference course, the intended interpretation is almost certainly the *causal* effect of R on X. Since the direct effect of R on X is positive in both strata of Z, the causal effect of rixq on xevu is positive.} \\
\\ 
\texttt{Thus, the correct answer is:} \\
\\ 
\texttt{Yes.} \\
            }
        }
    }
    & & \\ \\

    \theutterance \stepcounter{utterance}  
    & & & \multicolumn{2}{p{0.3\linewidth}}{
        \cellcolor[rgb]{0.9,0.9,0.9}{
            \makecell[{{p{\linewidth}}}]{
                \texttt{\tiny{[GM$|$GM]}}
                \texttt{No} \\
            }
        }
    }
    & & \\ \\

    \theutterance \stepcounter{utterance}  
    & & & \multicolumn{2}{p{0.3\linewidth}}{
        \cellcolor[rgb]{0.9,0.9,0.9}{
            \makecell[{{p{\linewidth}}}]{
                \texttt{\tiny{[GM$|$GM]}}
                \texttt{yes} \\
            }
        }
    }
    & & \\ \\

    \theutterance \stepcounter{utterance}  
    & & & \multicolumn{2}{p{0.3\linewidth}}{
        \cellcolor[rgb]{0.9,0.9,0.9}{
            \makecell[{{p{\linewidth}}}]{
                \texttt{\tiny{[GM$|$GM]}}
                \texttt{game\_result = LOSE} \\
            }
        }
    }
    & & \\ \\

\end{supertabular}
}

\end{document}
